\heading{数据结构与算法B-模拟题1-期末}
\noteinfo{题目来源：模拟题，参考解答由 AI 生成。本次整理提供 C 语言版与 Python 语言版；除程序代码语言外，两版题目内容保持一致。}

\textbf{一、选择题（每小题 3 分，共 30 分）}

\begin{question}
数据结构的三个基本要素是（\quad）。
\begin{enumerate}
  \item 顺序结构、链式结构、散列结构
  \item 逻辑结构、存储结构、算法
  \item 线性结构、树型结构、图状结构
  \item 以上都不是
\end{enumerate}
\begin{solution}
选 B。数据结构三要素为逻辑结构、存储结构和算法（数据操作）。
\end{solution}
\end{question}

\begin{question}
栈是一种后进先出表。若入栈的顺序为 $\{1,2,3,4\}$，则出栈的顺序不可能是（\quad）。
\begin{enumerate}
  \item $1234$
  \item $1243$
  \item $1423$
  \item $2134$
\end{enumerate}
\begin{solution}
选 C。$1423$ 不可能：要 $4$ 在 $2$ 之前出栈，则 $4$ 入栈时 $2,3$ 已在栈中，出栈顺序应为 $4,3,2$ 而非 $4,2,3$。
\end{solution}
\end{question}

\begin{question}
用一个大小为 $6$ 的数组来实现环形队列，元素从下标 $0$ 开始存放。当前 front 和 rear 的值分别为 $4$ 和 $1$，现在从队列中删除一个元素，再增加两个元素，则 front 和 rear 的值为（\quad）。
\begin{enumerate}
  \item $0,\;2$
  \item $5,\;3$
  \item $0,\;1$
  \item $5,\;2$
\end{enumerate}
\begin{solution}
选 B。环形队列删除：front = (4+1)\%6 = 5。增加两个：rear = (1+2)\%6 = 3。
\end{solution}
\end{question}

\begin{question}
假定一棵二叉树的节点个数为 $33$ 个，定义根节点的层数为 $0$，则它的高度最小为（\quad）。
\begin{enumerate}
  \item $5$
  \item $6$
  \item $7$
  \item $4$
\end{enumerate}
\begin{solution}
选 A。根结点层数为 $0$ 时，含 $33$ 个结点的二叉树最小最大层数为 $5$，因为前 $5$ 层最多有 $2^5-1=31$ 个结点，而前 $6$ 层最多有 $2^6-1=63$ 个结点。
\end{solution}
\end{question}

\begin{question}
已知某二叉树的先根周游序列为 $\{A,B,D,E,C\}$，中根周游序列为 $\{D,B,E,A,C\}$，则该二叉树的后根周游序列是（\quad）。
\begin{enumerate}
  \item $\{A,B,C,D,E\}$
  \item $\{D,B,E,C,A\}$
  \item $\{C,D,E,B,A\}$
  \item $\{D,E,B,C,A\}$
\end{enumerate}
\begin{solution}
选 D。由先序确定根 $A$，中序划分左子树 $\{D,B,E\}$、右子树 $\{C\}$。递归可得后序：$D,E,B,C,A$。
\end{solution}
\end{question}

\begin{question}
在待排序的元素序列基本有序的前提下，下列空间效率最差的排序方法是（\quad）。
\begin{enumerate}
  \item 直接插入排序
  \item 选择排序
  \item 快速排序
  \item 冒泡排序
\end{enumerate}
\begin{solution}
选 C。快速排序在基本有序时递归深度 $O(n)$，空间复杂度 $O(n)$，远高于其他 $O(1)$。
\end{solution}
\end{question}

\begin{question}
在平均情况下，下列时间复杂度最好的排序方法是（\quad）。
\begin{enumerate}
  \item 直接插入排序（平均 $O(n^2)$）
  \item 选择排序（平均 $O(n^2)$）
  \item 快速排序（平均 $O(n\log n)$）
  \item 冒泡排序（平均 $O(n^2)$）
\end{enumerate}
\begin{solution}
选 C。快速排序平均时间复杂度 $O(n\log n)$。
\end{solution}
\end{question}

\begin{question}
下列关于图的说法不正确的是（\quad）。
\begin{enumerate}
  \item 用邻接矩阵法存储一个图时，其存储空间与图的边数无关
  \item AOV 网络拓扑排序的结果是唯一的
  \item 强连通分量是有向图中的极大强连通子图
  \item 一个有向图的邻接表和逆邻接表中的结点个数一定相等
\end{enumerate}
\begin{solution}
选 B。拓扑排序结果不一定唯一，取决于入度为 $0$ 的顶点选择顺序。
\end{solution}
\end{question}

\begin{question}
AVL 树是一种（\quad）。
\begin{enumerate}
  \item 生成树
  \item 平衡树
  \item B 树
  \item 二叉排序树
\end{enumerate}
\begin{solution}
选 B。AVL 树是平衡二叉搜索树（任一节点左右子树高度差 $\le1$）。
\end{solution}
\end{question}

\begin{question}
任何一个无向连通图的最小生成树（\quad）。
\begin{enumerate}
  \item 只有一棵
  \item 有一棵或多棵
  \item 一定有多棵
  \item 可能不存在
\end{enumerate}
\begin{solution}
选 B。当图中有权值相等的边时，最小生成树可能不唯一。
\end{solution}
\end{question}

\clearpage

\textbf{二、解答题（共 40 分）}

\begin{question}[线性表存储结构比较]（8 分）

请说明线性表的顺序存储（顺序表）与链式存储（单链表）的异同，比较它们的优缺点。当字典采用二分法进行检索时，应采用哪一种存储结构？

\begin{solution}
\begin{itemize}
  \item 顺序表：逻辑上相邻的元素在物理上也相邻，支持随机访问 $O(1)$，插入/删除需移动元素 $O(n)$。
  \item 单链表：逻辑上相邻的元素物理上不一定相邻，通过指针链接，不支持随机访问 $O(n)$，插入/删除只需修改指针 $O(1)$。
  \item 二分法检索需要随机访问中间元素，应采用顺序表。
\end{itemize}
\end{solution}
\end{question}

\begin{question}[Huffman编码]（10 分）

字符串 \texttt{"this sentence is used for test"}，请使用 Huffman 编码方式对该字符串中的所有字符重新编码，以实现对它的压缩，给出编码过程中实现的 Huffman 树以及字符 \texttt{'s'} 的编码结果。

\begin{solution}
统计字符频率，构建 Huffman 树（每次合并频率最小的两个节点）。\texttt{'s'} 出现次数较多，其编码路径较短。具体 Huffman 树和编码需根据实际频率构建。
\end{solution}
\end{question}

\begin{question}[排序算法]（12 分）

对待排序的关键码集合 $\{E,C,A,M,S,R,D,F\}$ 按升序排序，请叙述直接插入排序、堆排序、快速排序、归并排序的基本思路，以及它们对上述数组第一趟扫描后的结果（快速排序使用第一个元素作为分界点）。

\begin{solution}
\begin{itemize}
  \item 直接插入排序：将第一个元素视为有序，依次将后续元素插入到有序部分的正确位置。第一趟后：$\{C,E,A,M,S,R,D,F\}$。
  \item 堆排序：先建堆，每次将堆顶与末尾元素交换并调整。第一趟后（建堆完成）：最大元素移至末尾。
  \item 快速排序：选 pivot（第一个元素 $E$），将小于 $E$ 的放左边、大于的放右边。第一趟后：$\{D,C,A,E,S,R,M,F\}$ 或类似。
  \item 归并排序：将序列两两归并。第一趟后：$\{C,E\},\{A,M\},\{R,S\},\{D,F\}$。
\end{itemize}
\end{solution}
\end{question}

\begin{question}[最短路径]（10 分）

下面是一个带权图邻接矩阵表示，试画出该图结构，并给出该图中由 $V_0$ 到 $V_4$ 的最短路径。

\begin{solution}
根据邻接矩阵画出带权有向图，使用 Dijkstra 算法求 $V_0$ 到 $V_4$ 的最短路径。
\end{solution}
\end{question}

\clearpage

\textbf{三、算法填空（共 10 分）}

\begin{question}[散列表检索]（10 分）

完成下列散列表的检索算法，假设散列函数为 $h$，采用线性探查法解决碰撞。

\begin{quote}
\footnotesize
\begin{tabular}{@{}l@{}}
typedef class \{\\
    DicElement element[REGION\_LEN];\\
    int m;  \# m=REGION\_LEN，为基本区域长度\\
\} HashDictionary;\\
\\
int linearSearch(HashDictionary *phash, KeyType key, int *position)\\
\{\\
    int d, inc;\\
    d = h(key);                               \# (1) 散列地址\\
    for (inc = 0; inc < phash.m; inc++) \{\\
        if (phash.element[d].key == key) \{   \# (2) 检索成功\\
            *position = d;\\
            return True;\\
        \}\\
        else if (phash.element[d].key == nil) \{\\
            *position = d;                    \# (3) 检索失败，找到插入位置\\
            return False;\\
        \}\\
        d = (d + 1) \% phash.m;              \# (4) 线性探查下一地址\\
    \}\\
    return False;                             \# (5) 散列表溢出\\
\}\\
\end{tabular}
\end{quote}

\begin{solution}
填空答案：(1) \texttt{h(key)} (2) \texttt{phash.element[d].key == key} (3) \texttt{*position = d} (4) \texttt{d = (d + 1) \% phash.m} (5) \texttt{return False}
\end{solution}
\end{question}

\clearpage

\textbf{四、算法设计（共 20 分）}

\begin{question}[求子树深度]（20 分）

已知一棵树采用下列存储定义，试求树中以值为 $x$ 的结点为根的子树深度。限定：值为 $x$ 的节点在树中最多有一个。注意栈或队列等辅助结构可以直接引用，不用自己定义。

\begin{quote}
\footnotesize
\begin{tabular}{@{}l@{}}
class CSNode;\\
typedef class CSNode *PCSNode;\\
class CSNode \{\\
    DataType info;\\
    PCSNode lchild;\\
    PCSNode rsibling;\\
\};\\
typedef class CSNode *CSTree;\\
CSTree pMytree;\\
\end{tabular}
\end{quote}

\begin{solution}
采用深度优先搜索（DFS）递归遍历：先找到值为 $x$ 的节点，然后递归计算其子树深度。
\begin{quote}
\footnotesize
\begin{tabular}{@{}l@{}}
int depth(CSTree root) \{\\
    if (root == None) return 0;\\
    int maxd = 0;\\
    CSTree child = root.lchild;\\
    while (child != None) \{\\
        int d = depth(child);\\
        if (d > maxd) maxd = d;\\
        child = child.rsibling;\\
    \}\\
    return maxd + 1;\\
\}\\
\end{tabular}
\end{quote}
先遍历全树找到值为 $x$ 的节点，再以其为根调用 \texttt{depth()}。
\end{solution}
\end{question}

\clearpage
